\chapter{Корневая стационарность общего Hodge--cycle профиля}
\label{ch:v7-common-carrier-root-stationarity-gate}

\section{Семь незамороженных корней}

Разрешим варьироваться трём старым заряженным рёбрам и четырём новым
singlet-рёбрам:
\begin{equation}
 r=(r_u,r_d,r_e,r_{LY},r_{XX},r_{Xe},r_{YY})\in\mathbb R_{>0}^7.
 \label{eq:v7-root-stationarity-roots}
\end{equation}
Старый сектор измеряется точным тепловым профилем хирального
Hodge-коммутатора, новый singlet-сектор --- профилем рёберного Hodge-момента,
а циклический сектор --- точным физическим Gaussian. Все три используют
одну шкалу $t=1$ и единичные представленные следы:
\begin{equation}
 \mathcal S_{\rm form}
 =-\Tr_{H_{15}}e^{-\mathfrak m_{\rm ch}^2}
  -\Tr_{E}e^{-\mathfrak m_C^2}
  +\Tr_{\rm phys}e^{-\Phi(r)^2}.
 \label{eq:v7-root-stationarity-formal-action}
\end{equation}

\section{Стационарная точка}

Совместные семь уравнений $\nabla_r\mathcal S_{\rm form}=0$ имеют решение
\begin{equation}
 r_*=(1.07737088,1.07737088,1.03492564,1.02768473,
      1.04972668,1.00214566,1.07390997).
 \label{eq:v7-root-stationarity-solution}
\end{equation}
Максимальный остаток градиента меньше $3\cdot10^{-15}$. Равенство первых
двух компонент следует сохранённой цвето-слабой симметрии старого
заряженного блока, а не наложено отдельно.

Корневой гессиан положителен:
\begin{equation}
 \Spec H_{rr}=\{15.61,16.92,18.01,18.90,20.65,
 60.22,60.22\}\quad\text{с округлением}.
 \label{eq:v7-root-stationarity-root-spectrum}
\end{equation}

\section{Полный смешанный гессиан}

Корневые и тяжёлые вариации не ортогональны: норма смешанного блока равна
\begin{equation}
 \|H_{r h}\|_{\max}=0.4780621494.
 \label{eq:v7-root-stationarity-cross-block}
\end{equation}
Поэтому отдельной положительности двух диагональных блоков недостаточно.
Прямое диагонализование полной матрицы размерности $7+20=27$ даёт
\begin{equation}
 (n_-,n_0,n_+)=(0,0,27),\qquad
 \lambda_{\min}=1.14399252085>0.
 \label{eq:v7-root-stationarity-full-signature}
\end{equation}

\begin{theorem}[Формальный локальный вакуумный проход]
При одной тепловой шкале и единичных следовых кратностях существует
положительная семикомпонентная корневая стационарная точка. Полный гессиан,
включающий все корневые и двадцать down--weak тяжёлых направлений, строго
положителен. Тем самым прежние tadpole и weak-препятствия одновременно
снимаются на уровне формального общего профиля.
\end{theorem}

\begin{nogocriterion}[Формальная сумма ещё не одна суперсвязность]
Результат нельзя повышать до физического родителя, пока три следа в
\eqref{eq:v7-root-stationarity-formal-action} не получены как один
физический полуслед кривизны одной Real-суперсвязности. Их знаки и
кратности пока согласованы минимально, но не выведены представленным
исчислением.
\end{nogocriterion}

\begin{protocol}
\begin{enumerate}
 \item незамороженных корневых амплитуд: \textbf{$7$};
 \item стационарное решение: \textbf{найдено};
 \item максимальный остаток градиента: \textbf{$<3\cdot10^{-15}$};
 \item положительных корневых мод: \textbf{$7$};
 \item тяжёлых мод: \textbf{$20$};
 \item полный смешанный гессиан: \textbf{$(0,0,27)$};
 \item минимальная полная мода: \textbf{$1.14399252085$};
 \item одна физическая Real-суперсвязность: \textbf{не выведена};
 \item итог: \textbf{формальный локальный проход};
 \item следующий гейт: \textbf{происхождение общего Real-полуследа}.
\end{enumerate}
\end{protocol}

Машинный сертификат сохранён в
\path{s2t_v7_common_carrier_root_stationarity_gate_results.json}.